% GENERATED FILE. Edit canonical lessons, then run npm run generate:books.
% canonical-source-hash: fnv1a64:75c9fbfa292106a8
% canonical-lessons: TE-C430-balaparacu, TE-C430-balahinaparacu, TE-C430-samatulyam, TE-C430-niyantrincu, TE-C430-parimiti

\chapter{Balance and Control}
\label{ch:te-balance-and-control}
\hlchaptermodality{430}

\begin{chapteropening}
\textbf{By the end of this chapter:} \emph{I can talk about strengthening, weakening, balance and limits.}

Complete the last lesson of chapter 430: I can talk about strengthening, weakening, balance and limits.
\end{chapteropening}

\section[\texorpdfstring{\te{బలపరచు} (balaparacu)}{బలపరచు (balaparacu)}]{\te{బలపరచు} (balaparacu) — to strengthen; to back up}

\label{lesson:TE-C430-balaparacu}



\begin{quote}
Before the new one: say the Telugu for to improve, then the Telugu for to turn into.
\end{quote}

\subsection*{You'll want to know: \te{బలపరచు}}
\textbf{\te{బలపరచు}} — \emph{balaparacu} — \textquotedblleft{}to strengthen; to back up\textquotedblright{}.

\begin{quote}
\te{ఆ} \te{సాక్ష్యం} \te{అతని} \te{వాదనను} \te{బలపరిచింది}.
\end{quote}

\begin{quote}
\emph{ā sākṣyaṁ atani vādananu balaparicindi.}
\end{quote}

\begin{quote}
That evidence strengthened his argument.
\end{quote}

Past (I …): \te{బలపరిచాను}, \emph{balaparicānu}.

\te{బలం}, strength, which you know, + -\te{పరచు}.

\begin{cousinweb}
\textbf{\te{బలపరచు}} (\emph{balaparacu}) is built from Sanskrit \emph{bala}, \textquotedblleft{}power, strength\textquotedblright{} (CDIAL 9161), the word in \te{బలం}, and -\te{పరచు}, a native ending that makes a verb (DEDR 3853 lists \emph{paṟacu}, \textquotedblleft{}to cause to feel\textquotedblright{}, beside \emph{paḍu}).
\end{cousinweb}

\begin{grammarlens}[title={What you've built}]
One more word for saying what happens if.
\end{grammarlens}

\subsection*{Your turn}
\begin{itemize}
\raggedright
  \item \emph{Say it:} \emph{balaparacu}
  \item \emph{Say it:} \emph{balaparacu}, once more
  \item \emph{Say it:} \emph{ā sākṣyaṁ atani vādananu balaparicindi.} — That evidence strengthened his argument.
  \item \emph{Recall:} say the Telugu for to improve, then the Telugu for to turn into, then say \emph{balaparacu} again
\end{itemize}

\begin{tcolorbox}[breakable,colback=teal!4,colframe=teal!35!black,title={Before you move on}]
What does \te{బలపరచు} mean? (\textquotedblleft{}To strengthen; to back up\textquotedblright{}.) Say it once more.
\end{tcolorbox}
\section[\texorpdfstring{\te{బలహీనపరచు} (balahīnaparacu)}{బలహీనపరచు (balahīnaparacu)}]{\te{బలహీనపరచు} (balahīnaparacu) — to weaken; to undermine}

\label{lesson:TE-C430-balahinaparacu}



\begin{quote}
Before the new one: say the Telugu for to turn into, then the Telugu for to strengthen.
\end{quote}

\subsection*{You'll want to know: \te{బలహీనపరచు}}
\textbf{\te{బలహీనపరచు}} — \emph{balahīnaparacu} — \textquotedblleft{}to weaken; to undermine\textquotedblright{}.

\begin{quote}
\te{అవినీతి} \te{ప్రజాస్వామ్యాన్ని} \te{బలహీనపరుస్తుంది}.
\end{quote}

\begin{quote}
\emph{avinīti prajāsvāmyānni balahīnaparustundi.}
\end{quote}

\begin{quote}
Corruption weakens democracy.
\end{quote}

Past (I …): \te{బలహీనపరిచాను}, \emph{balahīnaparicānu}.

\te{బలహీన}, weak, + -\te{పరచు}.

\begin{cousinweb}
\textbf{\te{బలహీనపరచు}} (\emph{balahīnaparacu}) is built from Sanskrit \emph{balahīna}, \textquotedblleft{}destitute of strength, weak\textquotedblright{} (Monier-Williams), and -\te{పరచు}, a native ending that makes a verb (DEDR 3853 lists \emph{paṟacu}, \textquotedblleft{}to cause to feel\textquotedblright{}, beside \emph{paḍu}).
\end{cousinweb}

\begin{grammarlens}[title={What you've built}]
One more word for saying what happens if.
\end{grammarlens}

\subsection*{Your turn}
\begin{itemize}
\raggedright
  \item \emph{Say it:} \emph{balahīnaparacu}
  \item \emph{Say it:} \emph{balahīnaparacu}, once more
  \item \emph{Say it:} \emph{avinīti prajāsvāmyānni balahīnaparustundi.} — Corruption weakens democracy.
  \item \emph{Recall:} say the Telugu for to turn into, then the Telugu for to strengthen, then say \emph{balahīnaparacu} again
\end{itemize}

\begin{tcolorbox}[breakable,colback=teal!4,colframe=teal!35!black,title={Before you move on}]
What does \te{బలహీనపరచు} mean? (\textquotedblleft{}To weaken; to undermine\textquotedblright{}.) Say it once more.
\end{tcolorbox}
\section[\texorpdfstring{\te{సమతుల్యం} (samatulyaṁ)}{సమతుల్యం (samatulyaṁ)}]{\te{సమతుల్యం} (samatulyaṁ) — balance; equilibrium}

\label{lesson:TE-C430-samatulyam}



\begin{quote}
Before the new one: say the Telugu for to strengthen, then the Telugu for to weaken.
\end{quote}

\subsection*{You'll want to know: \te{సమతుల్యం}}
\textbf{\te{సమతుల్యం}} — \emph{samatulyaṁ} — \textquotedblleft{}balance; equilibrium\textquotedblright{}.

\begin{quote}
\te{పనికి}, \te{కుటుంబానికి} \te{మధ్య} \te{సమతుల్యం} \te{ఉండాలి}.
\end{quote}

\begin{quote}
\emph{paniki, kuṭumbāniki madhya samatulyaṁ uṇḍāli.}
\end{quote}

\begin{quote}
There should be a balance between work and family.
\end{quote}

\te{సమతుల్య} \te{ఆహారం}, a balanced diet.

\begin{cousinweb}
\textbf{\te{సమతుల్యం}} (\emph{samatulyaṁ}) is built from Sanskrit \emph{sama}, \textquotedblleft{}same, even\textquotedblright{}, and \emph{tulya}, \textquotedblleft{}equal to, like\textquotedblright{} (CDIAL 5888).
\end{cousinweb}

\begin{grammarlens}[title={What you've built}]
One more word for saying what happens if.
\end{grammarlens}

\subsection*{Your turn}
\begin{itemize}
\raggedright
  \item \emph{Say it:} \emph{samatulyaṁ}
  \item \emph{Say it:} \emph{samatulyaṁ}, once more
  \item \emph{Say it:} \emph{paniki, kuṭumbāniki madhya samatulyaṁ uṇḍāli.} — There should be a balance between work and family.
  \item \emph{Recall:} say the Telugu for to strengthen, then the Telugu for to weaken, then say \emph{samatulyaṁ} again
\end{itemize}

\begin{tcolorbox}[breakable,colback=teal!4,colframe=teal!35!black,title={Before you move on}]
What does \te{సమతుల్యం} mean? (\textquotedblleft{}Balance; equilibrium\textquotedblright{}.) Say it once more.
\end{tcolorbox}
\section[\texorpdfstring{\te{నియంత్రించు} (niyantriñcu)}{నియంత్రించు (niyantriñcu)}]{\te{నియంత్రించు} (niyantriñcu) — to control; to regulate}

\label{lesson:TE-C430-niyantrincu}



\begin{quote}
Before the new one: say the Telugu for to weaken, then the Telugu for balance.
\end{quote}

\subsection*{You'll want to know: \te{నియంత్రించు}}
\textbf{\te{నియంత్రించు}} — \emph{niyantriñcu} — \textquotedblleft{}to control; to regulate\textquotedblright{}.

\begin{quote}
\te{ప్రభుత్వం} \te{ధరలను} \te{నియంత్రించాలి}.
\end{quote}

\begin{quote}
\emph{prabhutvaṁ dharalanu niyantriñcāli.}
\end{quote}

\begin{quote}
The government should control prices.
\end{quote}

Past (I …): \te{నియంత్రించాను}, \emph{niyantriñcānu}.

The noun is \te{నియంత్రణ}, control.

\begin{cousinweb}
\textbf{\te{నియంత్రించు}} (\emph{niyantriñcu}) is built on Sanskrit \emph{niyantraṇa}, \textquotedblleft{}restraining, checking, governing\textquotedblright{} (Monier-Williams), with the Telugu verb ending -\te{ఇంచు}.
\end{cousinweb}

\begin{grammarlens}[title={What you've built}]
One more word for saying what happens if.
\end{grammarlens}

\subsection*{Your turn}
\begin{itemize}
\raggedright
  \item \emph{Say it:} \emph{niyantriñcu}
  \item \emph{Say it:} \emph{niyantriñcu}, once more
  \item \emph{Say it:} \emph{prabhutvaṁ dharalanu niyantriñcāli.} — The government should control prices.
  \item \emph{Recall:} say the Telugu for to weaken, then the Telugu for balance, then say \emph{niyantriñcu} again
\end{itemize}

\begin{tcolorbox}[breakable,colback=teal!4,colframe=teal!35!black,title={Before you move on}]
What does \te{నియంత్రించు} mean? (\textquotedblleft{}To control; to regulate\textquotedblright{}.) Say it once more.
\end{tcolorbox}
\section[\texorpdfstring{\te{పరిమితి} (parimiti)}{పరిమితి (parimiti)}]{\te{పరిమితి} (parimiti) — a limit; a restriction}

\label{lesson:TE-C430-parimiti}



\begin{quote}
Before the new one: say the Telugu for balance, then the Telugu for to control.
\end{quote}

\subsection*{You'll want to know: \te{పరిమితి}}
\textbf{\te{పరిమితి}} — \emph{parimiti} — \textquotedblleft{}a limit; a restriction\textquotedblright{}.

\begin{quote}
\te{ప్రతి} \te{దానికీ} \te{ఒక} \te{పరిమితి} \te{ఉంటుంది}.
\end{quote}

\begin{quote}
\emph{prati dānikī oka parimiti uṇṭundi.}
\end{quote}

\begin{quote}
Everything has a limit.
\end{quote}

\te{పరిమితంగా}, in moderation.

\begin{cousinweb}
\textbf{\te{పరిమితి}} (\emph{parimiti}) is borrowed from Sanskrit \emph{parimiti}, \textquotedblleft{}measure, quantity, limitation\textquotedblright{} (Monier-Williams).
\end{cousinweb}

\begin{grammarlens}[title={What you've built}]
One more word for saying what happens if.
\end{grammarlens}

\subsection*{Your turn}
\begin{itemize}
\raggedright
  \item \emph{Say it:} \emph{parimiti}
  \item \emph{Say it:} \emph{parimiti}, once more
  \item \emph{Say it:} \emph{prati dānikī oka parimiti uṇṭundi.} — Everything has a limit.
  \item \emph{Recall:} say the Telugu for balance, then the Telugu for to control, then say \emph{parimiti} again
\end{itemize}

\begin{tcolorbox}[breakable,colback=teal!4,colframe=teal!35!black,title={Before you move on}]
What does \te{పరిమితి} mean? (\textquotedblleft{}A limit; a restriction\textquotedblright{}.) Say it once more.
\end{tcolorbox}
